\IfFileExists{stacks-project.cls}{%
\documentclass{stacks-project}
}{%
\documentclass{amsart}
}
\usepackage{amssymb}

% For dealing with references we use the comment environment
\usepackage{verbatim}
\usepackage{fancyvrb}
\usepackage{pmboxdraw}
\newenvironment{reference}{\comment}{\endcomment}
%\newenvironment{reference}{}{}
\newenvironment{slogan}{\comment}{\endcomment}
\newenvironment{history}{\comment}{\endcomment}

% For source figures retained as assets
\usepackage{graphicx}

% For commutative diagrams we use Xy-pic
\usepackage[all]{xy}

% We use 2cell for 2-commutative diagrams.
\xyoption{2cell}
\UseAllTwocells

% We use multicol for the list of chapters between chapters
\usepackage{multicol}

% This is generally recommended for better output
\usepackage{lmodern}
\usepackage[T1]{fontenc}
\usepackage{textalpha}

% For cross-file-references
% We prefer the snapshot tmp/xref/#2.aux,
% which latexmkrc replaces atomically, so
% that documents compiled simultaneously
% never read a half-written .aux file.
% The last argument keeps links to #2.pdf.
\usepackage{xr-hyper}
\newcommand{\maybeexternaldocument}[2]{%
  \IfFileExists{#2.aux}{%
    \externaldocument[#1]{#2}%
  }{%
    \PackageWarning{stack}{Missing #2.aux; external references with prefix #1 unavailable}%
  }%
}
\let\externaldocumentorig\externaldocument
\renewcommand{\externaldocument}[2][]{%
  \IfFileExists{tmp/xref/#2.aux}{%
    \externaldocumentorig[#1]{tmp/xref/#2}[#2.pdf]%
  }{%
  \IfFileExists{#2.aux}{%
    \externaldocumentorig[#1]{#2}%
  }{%
    \PackageWarning{stack}{Missing #2.aux; external references with prefix #1 unavailable}%
  }%
  }%
}

% Package for hypertext links:
\usepackage{hyperref}

% For any local file, say "hello.tex" you want to link to please
% use \externaldocument[hello-]{hello}
\externaldocument[introduction-]{introduction}
\externaldocument[conventions-]{conventions}
\externaldocument[sets-]{sets}
\externaldocument[categories-]{categories}
\externaldocument[topology-]{topology}
\externaldocument[sheaves-]{sheaves}
\externaldocument[sites-]{sites}
\externaldocument[stacks-]{stacks}
\externaldocument[fields-]{fields}
\externaldocument[algebra-]{algebra}
\externaldocument[brauer-]{brauer}
\externaldocument[homology-]{homology}
\externaldocument[derived-]{derived}
\externaldocument[simplicial-]{simplicial}
\externaldocument[more-algebra-]{more-algebra}
\externaldocument[smoothing-]{smoothing}
\externaldocument[modules-]{modules}
\externaldocument[sites-modules-]{sites-modules}
\externaldocument[injectives-]{injectives}
\externaldocument[cohomology-]{cohomology}
\externaldocument[sites-cohomology-]{sites-cohomology}
\externaldocument[dga-]{dga}
\externaldocument[dpa-]{dpa}
\externaldocument[sdga-]{sdga}
\externaldocument[hypercovering-]{hypercovering}
\externaldocument[schemes-]{schemes}
\externaldocument[constructions-]{constructions}
\externaldocument[properties-]{properties}
\externaldocument[morphisms-]{morphisms}
\externaldocument[coherent-]{coherent}
\externaldocument[divisors-]{divisors}
\externaldocument[limits-]{limits}
\externaldocument[varieties-]{varieties}
\externaldocument[topologies-]{topologies}
\externaldocument[descent-]{descent}
\externaldocument[perfect-]{perfect}
\externaldocument[more-morphisms-]{more-morphisms}
\externaldocument[flat-]{flat}
\externaldocument[groupoids-]{groupoids}
\externaldocument[more-groupoids-]{more-groupoids}
\externaldocument[etale-]{etale}
\externaldocument[chow-]{chow}
\externaldocument[intersection-]{intersection}
\externaldocument[pic-]{pic}
\externaldocument[weil-]{weil}
\externaldocument[adequate-]{adequate}
\externaldocument[dualizing-]{dualizing}
\externaldocument[duality-]{duality}
\externaldocument[discriminant-]{discriminant}
\externaldocument[derham-]{derham}
\externaldocument[local-cohomology-]{local-cohomology}
\externaldocument[algebraization-]{algebraization}
\externaldocument[curves-]{curves}
\externaldocument[resolve-]{resolve}
\externaldocument[models-]{models}
\externaldocument[functors-]{functors}
\externaldocument[equiv-]{equiv}
\externaldocument[pione-]{pione}
\externaldocument[etale-cohomology-]{etale-cohomology}
\externaldocument[proetale-]{proetale}
\externaldocument[relative-cycles-]{relative-cycles}
\externaldocument[more-etale-]{more-etale}
\externaldocument[trace-]{trace}
\externaldocument[crystalline-]{crystalline}
\externaldocument[spaces-]{spaces}
\externaldocument[spaces-properties-]{spaces-properties}
\externaldocument[spaces-morphisms-]{spaces-morphisms}
\externaldocument[decent-spaces-]{decent-spaces}
\externaldocument[spaces-cohomology-]{spaces-cohomology}
\externaldocument[spaces-limits-]{spaces-limits}
\externaldocument[spaces-divisors-]{spaces-divisors}
\externaldocument[spaces-over-fields-]{spaces-over-fields}
\externaldocument[spaces-topologies-]{spaces-topologies}
\externaldocument[spaces-descent-]{spaces-descent}
\externaldocument[spaces-perfect-]{spaces-perfect}
\externaldocument[spaces-more-morphisms-]{spaces-more-morphisms}
\externaldocument[spaces-flat-]{spaces-flat}
\externaldocument[spaces-groupoids-]{spaces-groupoids}
\externaldocument[spaces-more-groupoids-]{spaces-more-groupoids}
\externaldocument[bootstrap-]{bootstrap}
\externaldocument[spaces-pushouts-]{spaces-pushouts}
\externaldocument[spaces-chow-]{spaces-chow}
\externaldocument[groupoids-quotients-]{groupoids-quotients}
\externaldocument[spaces-more-cohomology-]{spaces-more-cohomology}
\externaldocument[spaces-simplicial-]{spaces-simplicial}
\externaldocument[spaces-duality-]{spaces-duality}
\externaldocument[formal-spaces-]{formal-spaces}
\externaldocument[restricted-]{restricted}
\externaldocument[spaces-resolve-]{spaces-resolve}
\externaldocument[formal-defos-]{formal-defos}
\externaldocument[defos-]{defos}
\externaldocument[cotangent-]{cotangent}
\externaldocument[examples-defos-]{examples-defos}
\externaldocument[algebraic-]{algebraic}
\externaldocument[examples-stacks-]{examples-stacks}
\externaldocument[stacks-sheaves-]{stacks-sheaves}
\externaldocument[criteria-]{criteria}
\externaldocument[artin-]{artin}
\externaldocument[quot-]{quot}
\externaldocument[stacks-properties-]{stacks-properties}
\externaldocument[stacks-morphisms-]{stacks-morphisms}
\externaldocument[stacks-limits-]{stacks-limits}
\externaldocument[stacks-cohomology-]{stacks-cohomology}
\externaldocument[stacks-perfect-]{stacks-perfect}
\externaldocument[stacks-introduction-]{stacks-introduction}
\externaldocument[stacks-more-morphisms-]{stacks-more-morphisms}
\externaldocument[stacks-geometry-]{stacks-geometry}
\externaldocument[moduli-]{moduli}
\externaldocument[moduli-curves-]{moduli-curves}
\externaldocument[examples-]{examples}
\externaldocument[exercises-]{exercises}
\externaldocument[guide-]{guide}
\externaldocument[desirables-]{desirables}
\externaldocument[coding-]{coding}
\externaldocument[obsolete-]{obsolete}
\externaldocument[fdl-]{fdl}
\externaldocument[index-]{index}

% Heap Project documents
\externaldocument[k3-]{k3}
\externaldocument[torelli-]{torelli}
\externaldocument[linear-algebra-]{linear-algebra}
\externaldocument[groups-]{groups}
\externaldocument[complex-analysis-]{complex-analysis}
\externaldocument[functional-analysis-]{functional-analysis}
\externaldocument[categories-]{categories}
\externaldocument[commutative-algebra-]{commutative-algebra}
\externaldocument[ringed-spaces-]{ringed-spaces}
\externaldocument[homological-algebra-]{homological-algebra}
\externaldocument[lie-algebras-]{lie-algebras}
\externaldocument[representation-theory-]{representation-theory}
\externaldocument[smooth-manifolds-]{smooth-manifolds}
\externaldocument[differential-topology-]{differential-topology}
\externaldocument[algebraic-topology-]{algebraic-topology}
\externaldocument[differential-geometry-]{differential-geometry}
\externaldocument[algebraic-geometry-]{algebraic-geometry}
\externaldocument[symplectic-geometry-]{symplectic-geometry}
\externaldocument[complex-geometry-]{complex-geometry}
\externaldocument[hodge-theory-]{hodge-theory}
\externaldocument[%
variations-of-hodge-structures-]{%
variations-of-hodge-structures}
\externaldocument[%
homological-mirror-symmetry-]{%
homological-mirror-symmetry}
\maybeexternaldocument{deformations-}{deformations}
\externaldocument[vertex-algebras-]{vertex-algebras}
\externaldocument[springer-]{springer}
\externaldocument[infty-categories-]{infty-categories}
\externaldocument[seminars-differential-geometry-]{%
seminars-differential-geometry}
\externaldocument[seminars-complex-geometry-]{%
seminars-complex-geometry}
\externaldocument[seminars-algebraic-geometry-]{%
seminars-algebraic-geometry}
\externaldocument[seminars-analysis-pde-]{%
seminars-analysis-pde}
\externaldocument[surfaces-]{surfaces}
\externaldocument[birational-maps-commutative-algebra-]{birational-maps-commutative-algebra}
\externaldocument[cremona-transformations-]{cremona-transformations}
\externaldocument[derived-categories-of-sheaves-]{derived-categories-of-sheaves}
\externaldocument[geometric-invariant-theory-]{geometric-invariant-theory}
\externaldocument[geometric-stability-conditions-and-group-actions-]{geometric-stability-conditions-and-group-actions}
\externaldocument[moduli-spaces-of-sheaves-]{moduli-spaces-of-sheaves}
\externaldocument[bridgeland-stability-general-theory-]{bridgeland-stability-general-theory}
\externaldocument[hilbert-schemes-]{hilbert-schemes}
\externaldocument[noncommutative-abelian-surfaces-and-kummer-type-hyperkahler-varieties-]{noncommutative-abelian-surfaces-and-kummer-type-hyperkahler-varieties}
\externaldocument[mumford-tate-groups-in-hodge-theory-]{mumford-tate-groups-in-hodge-theory}
\externaldocument[hodge-birational-atoms-2026-]{hodge-birational-atoms-2026}
\externaldocument[xxii-egd-2026-]{%
xxii-egd-2026}
\externaldocument[pde-]{pde}
\externaldocument[probability-]{probability}

\externaldocument[linguistics-]{linguistics}

\externaldocument[%
16th-alga-meeting-2026-]{%
16th-alga-meeting-2026}

\externaldocument[%
iv-jornadas-lefschetz-2026-]{%
iv-jornadas-lefschetz-2026}

\externaldocument[reading-the-masters-]{%
reading-the-masters}

\externaldocument[real-analysis-]{%
real-analysis}

\externaldocument[optimal-transport-]{%
optimal-transport}

\externaldocument[history-]{history}

\externaldocument[seiberg-witten-]{seiberg-witten}

\externaldocument[history-philosophy-biology-]{history-philosophy-biology}

\externaldocument[classical-mechanics-]{classical-mechanics}
\externaldocument[electromagnetism-]{electromagnetism}
\externaldocument[waves-and-optics-]{waves-and-optics}
\externaldocument[wif-]{waves-in-fluids-iv}
\externaldocument[thermal-physics-]{thermal-physics}
\externaldocument[quantum-mechanics-]{quantum-mechanics}
\externaldocument[special-relativity-]{special-relativity}
\externaldocument[physics-]{physics}

\externaldocument[chemistry-]{chemistry}

\externaldocument[molecular-cell-biology-]{molecular-cell-biology}

\externaldocument[microbiology-]{microbiology}

\externaldocument[protistology-]{protistology}

\externaldocument[mycology-]{mycology}

\externaldocument[botany-]{botany}

\externaldocument[zoology-]{zoology}

\externaldocument[animal-histology-]{animal-histology}

\externaldocument[animal-physiology-]{animal-physiology}

\externaldocument[systematics-biogeography-]{systematics-biogeography}

\externaldocument[ecology-natural-resources-]{ecology-natural-resources}

\externaldocument[earth-sciences-]{earth-sciences}

\externaldocument[biostatistics-]{biostatistics}

\externaldocument[mathematical-systems-biology-]{mathematical-systems-biology}

\externaldocument[cimpa-el-salvador-2024-]{cimpa-el-salvador-2024}

\externaldocument[anthropology-]{anthropology}

\externaldocument[yang-mills-theory-]{yang-mills-theory}

\externaldocument[seminars-algebra-]{seminars-algebra}

\externaldocument[seminars-physics-]{seminars-physics}

\externaldocument[seminars-computing-]{seminars-computing}

\externaldocument[oktobermat-2026-]{oktobermat-2026}

\externaldocument[saet-2026-]{saet-2026}

\externaldocument[seminars-symplectic-geometry-]{seminars-symplectic-geometry}

\externaldocument[seminars-dynamics-]{seminars-dynamics}

\externaldocument[oktobermat-2024-]{oktobermat-2024}

\externaldocument[quantum-groups-]{quantum-groups}

\externaldocument[legal-2027-]{legal-2027}

\externaldocument[derived-categories-]{derived-categories}

% Theorem environments.
%
\theoremstyle{plain}
\newtheorem{theorem}[subsection]{Theorem}
\newtheorem{proposition}[subsection]{Proposition}
\newtheorem{lemma}[subsection]{Lemma}
\newtheorem{corollary}[subsection]{Corollary}

\theoremstyle{definition}
\newtheorem{definition}[subsection]{Definition}
\newtheorem{example}[subsection]{Example}
\newtheorem{exercise}[subsection]{Exercise}
\newtheorem{situation}[subsection]{Situation}

\theoremstyle{remark}
\newtheorem{remark}[subsection]{Remark}
\newtheorem{remarks}[subsection]{Remarks}

\numberwithin{equation}{subsection}

% Macros
%
\def\lim{\mathop{\mathrm{lim}}\nolimits}
\def\colim{\mathop{\mathrm{colim}}\nolimits}
\def\Spec{\mathop{\mathrm{Spec}}}
\def\Hom{\mathop{\mathrm{Hom}}\nolimits}
\def\Ext{\mathop{\mathrm{Ext}}\nolimits}
\def\SheafHom{\mathop{\mathcal{H}\!\mathit{om}}\nolimits}
\def\SheafExt{\mathop{\mathcal{E}\!\mathit{xt}}\nolimits}
\def\Sch{\mathit{Sch}}
\def\Mor{\mathop{\mathrm{Mor}}\nolimits}
\def\Ob{\mathop{\mathrm{Ob}}\nolimits}
\def\Sh{\mathop{\mathit{Sh}}\nolimits}
\def\NL{\mathop{N\!L}\nolimits}
\def\CH{\mathop{\mathrm{CH}}\nolimits}
\def\proetale{{pro\text{-}\acute{e}tale}}
\def\etale{{\acute{e}tale}}
\def\QCoh{\mathit{QCoh}}
\def\Ker{\mathop{\mathrm{Ker}}}
\def\Im{\mathop{\mathrm{Im}}}
\def\Coker{\mathop{\mathrm{Coker}}}
\def\Coim{\mathop{\mathrm{Coim}}}

% Boxtimes
%
\DeclareMathSymbol{\boxtimes}{\mathbin}{AMSa}{"02}

%
% Macros for moduli stacks/spaces
%
\def\QCohstack{\mathcal{QC}\!\mathit{oh}}
\def\Cohstack{\mathcal{C}\!\mathit{oh}}
\def\Spacesstack{\mathcal{S}\!\mathit{paces}}
\def\Quotfunctor{\mathrm{Quot}}
\def\Hilbfunctor{\mathrm{Hilb}}
\def\Curvesstack{\mathcal{C}\!\mathit{urves}}
\def\Polarizedstack{\mathcal{P}\!\mathit{olarized}}
\def\Complexesstack{\mathcal{C}\!\mathit{omplexes}}
% \Pic is the operator that assigns to X its picard group, usage \Pic(X)
% \Picardstack_{X/B} denotes the Picard stack of X over B
% \Picardfunctor_{X/B} denotes the Picard functor of X over B
\def\Pic{\mathop{\mathrm{Pic}}\nolimits}
\def\Picardstack{\mathcal{P}\!\mathit{ic}}
\def\Picardfunctor{\mathrm{Pic}}
\def\Deformationcategory{\mathcal{D}\!\mathit{ef}}


\begin{document}

\title{Seminars in Algebra}
\maketitle

\phantomsection
\label{section-phantom}

\tableofcontents

\section{Classification of admissible
W-algebras of low growth}
\label{section-w-algebras-of-low-growth}

\noindent
Igor Alarcon Blatt. June 10, 2026.
Algebra Seminar, IMPA.

\medskip\noindent
{\bf Abstract.} 
Admissible W-algebras form an important class
of vertex algebras, closely related to affine
Kac-Moody algebras and conformal field
theory. A natural invariant of such algebras
is (the asymptotic behaviour of) their
normalised characters. We classify all
admissible W-algebras whose normalised
characters have asymptotic growth less than
1. The proof relies on the computation of the
relevant data together with a theorem of Kac
and Wakimoto, leading to a characterisation
in terms of Virasoro minimal models.

\medskip\noindent
Goal: classify a certain class of W-algebras.
W-algebras are a kind of vertex algebra.

[A lot of content missing.
Some ideas are: minimal models
of Virasoro algebras,
assymptotic data]

\medskip\noindent
Recall that partitions of $n$
correspond with nilpotent orbits.

\begin{theorem}[Arakawa, 2015]
\label{theorem-arakawa}
For $k$ an admissible level with
denominator $u$, there is a unique
nilpotent orbit $\mathbb{O}_n$
suh that $W_k(\mathfrak{g},f)$
is rational for $f \in \overline{
\mathbb{O}_n}$. The orbits $\mathbb{O}_n$
correspond to the partitions listed below.
\end{theorem}

\medskip\noindent
Upshot: use a lemma 
(by Kac-Wakimoto, perhaps in
``Quantum reduction
for affine superalgebras'') which asserts
that if the assymptotic growth
satisfies $>1$ then there is a direct
sum expression.



\section{Vertex Algebras and the Monster}
\label{section-vertex-algebras-monster}

\noindent
David Melo. September 18, 2026.
Students' Colloquium. Room 228.
Friday, 15:20--16:20.

\medskip\noindent
{\bf Abstract.}
The classification of finite simple groups is
one of the most remarkable achievements in
algebra. Its proof spans about 500 articles,
involves about 100 mathematicians, and fills
more than 10,000 pages. The work surrounding
the classification also uncovered exceptional
structures connecting several areas of
mathematics.

A particularly striking example is the
relationship among the Monster group, sphere
packings, modular forms, and conformal field
theory through vertex algebras. The Leech
lattice gives the optimal sphere packing in
dimension $24$. A conformal field theory is
constructed from this lattice. Borcherds's
proof of monstrous moonshine explains
remarkable coincidences between the
Monster and Ramanujan's $j$-function. This
talk will introduce the exceptional objects
and outline their relationship.

\medskip\noindent
{\bf Introduction.}
The Monster is the largest sporadic finite
simple group. Its order is
$$
\begin{aligned}
|\mathbb{M}|
&=2^{46}3^{20}5^9 7^6 11^2 13^3 17\\
&\phantom{={}}\cdot19\cdot23\cdot29\cdot31
\cdot41\cdot47\cdot59\cdot71.
\end{aligned}
$$
Some of the coefficients in the Fourier
expansion of Ramanujan's $j$-function have
unexpected connections with dimensions of
representations of the Monster. Monstrous
moonshine explains these numerical
coincidences.

The Monster is the automorphism group of the
moonshine vertex algebra.

\medskip\noindent
{\bf Applications of the modular group.}
The modular group has applications in the
following areas.

\begin{enumerate}
\item
It parametrizes complex structures on a
torus.

\item
It occurs in Diophantine equations, including
the proof of Fermat's last theorem.

\item
It is used in cryptography.

\item
It helps compute elliptic integrals.

\item
It governs Fourier expansions of modular
functions and modular forms.
\end{enumerate}

\medskip\noindent
{\bf Modular invariants.}
The modular group is
$$
\operatorname{SL}_2(\mathbb{Z}).
$$
Its action on the upper half-plane comes from
translation and inversion.
$$
\tau\longmapsto\tau+1,
\qquad
\tau\longmapsto-\frac{1}{\tau}.
$$

A modular function is a meromorphic invariant
function for this action. A modular form of
weight $k$ transforms by
$$
f\left(\frac{a\tau+b}{c\tau+d}\right)
=(c\tau+d)^k f(\tau).
$$
The matrix with entries $a,b,c,d$ belongs
to the modular group.

Examples include Eisenstein series, theta
functions, and the Dedekind eta function.
Their coefficients often count solutions of
arithmetic equations. They were a key
ingredient in the proof of Fermat's last
theorem.

The $j$-function is the simplest nonconstant
modular function. Its Fourier expansion is
$$
j(\tau)=q^{-1}+744+196884q
+21493760q^2+\cdots.
$$

\medskip\noindent
{\bf Classification of simple groups.}
Let $G$ be a finite simple group. Then $G$ is
one of the following.

\begin{enumerate}
\item
A cyclic group of prime order.

\item
An alternating group.

\item
A group of Lie type.

\item
One of the $26$ sporadic groups.
\end{enumerate}

The periodic table of finite simple groups
displays the sporadic groups and their
relationships.

$$
\xymatrix{
& \mathbb{M}\ar[dl]\ar[d]\ar[dr] & \\
\mathbb{B} &
\mathrm{Fi}_{24}' &
\mathrm{Co}_1 \\
& \mathrm{Fi}_{23}\ar[ur] &
\mathrm{Co}_2\ar[u]
}
$$

\medskip\noindent
{\bf Lattices.}
An even lattice is a discrete subgroup
$L\subset\mathbb{R}^n$ such that
$$
\left\langle{} x,x \right\rangle{}
\in2\mathbb{Z}
$$
for every $x\in L$. Its dual lattice consists
of vectors in $\mathbb{R}^n$ whose pairings
with every element of $L$ are integral.

The roots of an even lattice are the vectors
of norm $2$. They have the smallest possible
positive norm.

\medskip\noindent
{\bf The Leech lattice.}
The Leech lattice comes from coding theory.
It gives the optimal sphere packing in
dimension $24$, as was proved in 2016. It is
the unique even unimodular lattice of rank
$24$ without roots.

The identity
$$
1^2+2^2+\cdots+24^2=70^2
$$
is another striking occurrence of the number
$24$.

Show a slide with a matrix whose rows form a
basis of the Leech lattice.

\medskip\noindent
{\bf Conway's group.}
The Conway group is
$$
\mathrm{Co}_1=
\operatorname{Aut}(\Lambda)/\{1,-1\}.
$$
It is a sporadic simple group. The Monster is
not obtained from $\mathrm{Co}_1$ by
adjoining one element. Rather,
$\mathrm{Co}_1$ occurs as a quotient of a
subgroup of the Monster.

\medskip\noindent
{\bf Group representations.}
{\it Definition.} A representation of $G$ on
$V$ is a homomorphism
$$
\rho:G\longrightarrow\operatorname{GL}(V).
$$
A morphism of representations is a linear map
$f:V\longrightarrow W$ such that
$$
f\rho(g)=\sigma(g)f
$$
for every $g\in G$.

Examples include the following.

\begin{enumerate}
\item
The group $\operatorname{Aut}(\Lambda)$ acts
on $\mathbb{R}^{24}$. The quotient
$\mathrm{Co}_1$ acts projectively.

\item
The Monster has an irreducible representation
of dimension $196883$.

\item
The symmetric group $S_n$ acts on
$\mathbb{R}^n$ by permuting coordinates.

\item
The circle has one-dimensional
representations indexed by $n\in\mathbb{Z}$.
$$
e^{i\theta}\longmapsto e^{in\theta}.
$$

\item
The groups $A_4$ and $A_5$ act on
$\mathbb{R}^3$ as rotation groups.
\end{enumerate}

\medskip\noindent
{\bf Putting a ring on it.}
For a group $G$, the group algebra
$\mathbb{k}[G]$ is the free vector space on
$G$, with multiplication extending the group
operation. Representations of $G$ are
equivalent to $\mathbb{k}[G]$-modules.

\medskip\noindent
{\bf Reducibility.}
{\it Definition.} A representation is
reducible if it has a nonzero proper
subrepresentation. In a semisimple category,
it is then a direct sum of two nonzero
representations. Otherwise, it is
irreducible.

\begin{theorem}
\label{theorem-finite-group-representations}
Let $G$ be finite and let $\mathbb{k}$ have
characteristic zero. There are finitely many
irreducible representations $V_j$, and every
finite-dimensional representation is a direct
sum of such representations.
\end{theorem}

\begin{theorem}[Peter--Weyl]
\label{theorem-peter-weyl}
Let $G$ be a compact topological group. Then
$$
L^2(G)=\bigoplus_{j\in\Sigma}
V_j\otimes V_j^*
$$
as an orthogonal decomposition of Hilbert
spaces.
\end{theorem}

For $G=S^1$, this gives
$$
L^2(S^1)=\bigoplus_{n\in\mathbb{Z}}
\mathbb{C}e^{in\theta}.
$$
The number of irreducible complex
representations of a finite group equals its
number of conjugacy classes.

\medskip\noindent
{\bf The Virasoro Lie algebra.}
{\it Definition.} A Lie algebra is a vector
space with a bilinear bracket satisfying
antisymmetry and the Jacobi identity.

On $\mathbb{C}^{\times}$, put
$$
L_n=-z^{n+1}\frac{d}{dz}.
$$
These vector fields satisfy
$$
[L_m,L_n]=(m-n)L_{m+n}.
$$
Their unique nontrivial central extension has
relations
$$
[L_m,L_n]=(m-n)L_{m+n}
+\frac{m^3-m}{12}\delta_{m+n,0}C.
$$
This is the Virasoro algebra. It is the
quantum analogue of local conformal
transformations and is central in string
theory. The central extension accounts for
normal-ordering anomalies.

\medskip\noindent
{\bf Positive-energy representations.}
{\it Definition.} A positive-energy
representation of the Virasoro algebra is a
vector space $V$. The element $C$ acts by a
scalar $c$, $L_0$ acts diagonally with
spectrum bounded below, and every eigenspace
of $L_0$ has
finite dimension.

Its renormalized character is
$$
\chi_V(q)=\sum_{h\in\mathbb{R}}
\dim(V^{L_0=h})q^{h-c/24}.
$$

\medskip\noindent
{\bf Vertex algebras.}
{\it Definition.} A conformal vertex algebra
consists of a vertex algebra $V$, a vacuum
vector $|0\rangle$, and a Virasoro vector $L$
with
$$
T=L_{-1}.
$$
The vector $L$ is the energy-momentum tensor.

\medskip\noindent
{\bf Borcherds's identity.}
For $a,b,c\in V$ and $m,n,k\in\mathbb{Z}$,
the identity is
$$
\begin{aligned}
&\sum_{j\geq0}\binom{m}{j}
(a_{(n+j)}b)_{(m+k-j)}c\\
&=\sum_{j\geq0}(-1)^j\binom{n}{j}
\left(a_{(m+n-j)}b_{(k+j)}c\right.\\
&\qquad\left.{}-(-1)^n
b_{(n+k-j)}a_{(m+j)}c\right).
\end{aligned}
$$

\medskip\noindent
{\bf Examples and applications.}
For a simple Lie algebra $\mathfrak{g}$ and a
complex number $k$, the universal affine
vertex algebra is $V^k(\mathfrak{g})$. Its
weight-one subspace is naturally
$\mathfrak{g}$. If $f\in\mathfrak{g}$ is
nilpotent, quantum Hamiltonian reduction
gives the vertex algebra
$W^k(\mathfrak{g},f)$. An even lattice $L$
gives a lattice vertex algebra
$V_L$.

Vertex algebras have applications to knot
theory and braid relations, integrable
partial differential equations in
symplectic geometry,
$D$-modules on Riemann surfaces,
combinatorial product-sum identities, and
theoretical physics.

\medskip\noindent
{\bf Bosonic strings.}
Let
$$
\mathcal{H}=\mathbb{C}[x_{-1},x_{-2},\ldots]
$$
and let $|0\rangle=1$. The Heisenberg modes
act by
$$
\begin{aligned}
(x_{-m})_{(n)}P&=x_{-m}P\qquad(n=-m),\\
(x_{-m})_{(n)}P&=n
\frac{\partial P}{\partial x_{-n}}
\qquad(n>0).
\end{aligned}
$$
For the standard conformal vector, the
central charge is $1$ and
$$
L_0x_{-k}=kx_{-k}.
$$
The character is
$$
\chi_{\mathcal{H}}(q)=
\prod_{m\geq1}\frac{1}{1-q^m}.
$$

\medskip\noindent
{\bf Lattice vertex algebras.}
For an even lattice $L$, the lattice vertex
algebra $V_L$ combines the Heisenberg algebra
with a copy of the group algebra of $L$.
Show a picture of the lattice, with a
Heisenberg cone at each lattice point.

\medskip\noindent
{\bf Moonshine.}
Starting with the Leech lattice vertex
algebra $V_\Lambda$, orbifolding by the
involution
induced by $-1$ gives the moonshine module
$$
V^{\natural}=V_\Lambda^+
\oplus V_\Lambda^{T,+}.
$$
The second summand comes from a twisted
module.
The automorphism group of $V^{\natural}$ is
the Monster.

\medskip\noindent
{\bf Extra ingredients.}
The no-ghost theorem states that critical
bosonic string theory in dimension $26$ has
no negative-norm physical states. Borcherds's
proof constructs a Lie algebra from a
semi-infinite cohomology group for
$$
V^{\natural}\otimes V_{II_{1,1}}.
$$
This construction is closely related to
current research on semi-infinite cohomology.

\bibliography{my}
\bibliographystyle{amsalpha}

\end{document}
